\newcounter{stoppagebreaks}
\documentclass{article}

% packages
\usepackage{amsmath, amsthm, thmtools, amsfonts, amssymb, luacode, catchfile, hyperref, ifthen, graphicx, bbm}
\ifcsname c@kobocompile\endcsname
	\usepackage[a5paper, total={1072pt, 1448pt}, margin=10pt, includeheadfoot]{geometry} % set page margins
\else
	\usepackage[a4paper, margin=50pt, includeheadfoot]{geometry}
\fi
\usepackage[shortlabels]{enumitem}
\usepackage[skip=3pt, indent=0pt]{parskip}

% language
\usepackage[bidi=basic, layout=tabular, provide=*]{babel}
\ifcsname c@english\endcsname
	\babelprovide[main, import]{english}
\else
	\babelprovide[main, import, alph=letters.gershayim]{hebrew}
	\setlist[enumerate,2]{label={\alph*.}}
	\babelprovide{rl}
\fi
\babelfont{rm}[Renderer=Harfbuzz]{Libertinus Serif}
\babelfont{sf}{Libertinus Sans}
\babelfont{tt}{Libertinus Mono}

% style
\ifcsname c@stoppagebreaks\endcsname
\else
	\AddToHook{cmd/section/before}{\clearpage}	% Add line break before section
\fi
\linespread{1.3}
\setcounter{secnumdepth}{0}		% Remove default number tags from sections, this won't do well with theorems
\AtBeginDocument{\setlength{\belowdisplayskip}{3pt}}
\AtBeginDocument{\setlength{\abovedisplayskip}{3pt}}
\graphicspath{ {../images/} }

% operators
\DeclareMathOperator\cis{cis}
\DeclareMathOperator\Sp{Span}
\DeclareMathOperator\tr{tr}
\DeclareMathOperator\im{Im}
\DeclareMathOperator\re{Re}
\DeclareMathOperator\diag{diag}
\DeclareMathOperator*\lowlim{\underline{lim}}
\DeclareMathOperator*\uplim{\overline{lim}}
\DeclareMathOperator\rng{rng}
\DeclareMathOperator\Sym{Sym}
\DeclareMathOperator\Arg{Arg}
\DeclareMathOperator\Log{Log}
\DeclareMathOperator\dom{dom}
\DeclareMathOperator\supp{Supp}
\DeclareMathOperator\var{Var}
\DeclareMathOperator\cov{Cov}
\DeclareMathOperator\End{End}
\DeclareMathOperator\id{id}
\DeclareMathOperator\cl{cl}
\DeclareMathOperator\gal{Gal}
\DeclareMathOperator\dist{dist}
\DeclareMathOperator\form{form}
\DeclareMathOperator\term{term}
\DeclareMathOperator\sent{sent}
\DeclareMathOperator\Th{Th}
\DeclareMathOperator\Hom{Hom}
\DeclareMathOperator\Aut{Aut}
\DeclareMathOperator\Char{char}
\DeclareMathOperator\vol{vol}
\DeclareMathOperator\GL{GL}
\DeclareMathOperator\SL{SL}
\DeclareMathOperator\lcm{lcm}

% commands
\newcommand{\cupdot}{\mathbin{\mathaccent\cdot\cup}}
\newcommand{\indicator}{\mathbbm{1}}
\newcommand{\AAA}[0]{\mathbb{A}}
\newcommand{\Aa}[0]{\mathcal{A}}
\newcommand{\BB}[0]{\mathbb{B}}
\newcommand{\Bb}[0]{\mathcal{B}}
\newcommand{\CC}[0]{\mathbb{C}}
\newcommand{\Cc}[0]{\mathcal{C}}
\newcommand{\DD}[0]{\mathbb{D}}
\newcommand{\Dd}[0]{\mathcal{D}}
\newcommand{\EE}[0]{\mathbb{E}}
\newcommand{\FF}[0]{\mathbb{F}}
\newcommand{\Ff}[0]{\mathcal{F}}
\newcommand{\Ii}[0]{\mathcal{I}}
\newcommand{\Gg}[0]{\mathcal{G}}
\newcommand{\HH}[0]{\mathbb{H}}
\newcommand{\Hh}[0]{\mathcal{H}}
\newcommand{\Kk}[0]{\mathcal{K}}
\newcommand{\Ll}[0]{\mathcal{L}}
\newcommand{\Mm}[0]{\mathcal{M}}
\newcommand{\NN}[0]{\mathbb{N}}
\newcommand{\Nn}[0]{\mathcal{N}}
\newcommand{\PP}[0]{\mathbb{P}}
\newcommand{\Pp}[0]{\mathcal{P}}
\newcommand{\QQ}[0]{\mathbb{Q}}
\newcommand{\RR}[0]{\mathbb{R}}
\newcommand{\Rr}[0]{\mathcal{R}}
\newcommand{\Ss}[0]{\mathcal{S}}
\newcommand{\TT}[0]{\mathbb{T}}
\newcommand{\Tt}[0]{\mathcal{T}}
\newcommand{\UU}[0]{\mathbb{U}}
\newcommand{\Uu}[0]{\mathcal{U}}
\newcommand{\Vv}[0]{\mathcal{V}}
\newcommand{\Ww}[0]{\mathcal{W}}
\newcommand{\Xx}[0]{\mathcal{X}}
\newcommand{\Yy}[0]{\mathcal{Y}}
\newcommand{\ZZ}[0]{\mathbb{Z}}
\newcommand{\Zz}[0]{\mathcal{Z}}
\newcommand{\acts}[0]{\circlearrowright}
\newcommand{\explain}[2] {
	\begin{flalign*}
		 && \text{#2} && \text{#1}
	\end{flalign*}
}
\newcommand{\maketitleprint}[1][green]{
	\begin{center}
		\includegraphics[width=6cm]{dragon_#1}
	\end{center}
}

% theorem commands
\newtheoremstyle{c_remark}
	{}	% Space above
	{}	% Space below
	{}% Body font
	{}	% Indent amount
	{\bfseries}	% Theorem head font
	{}	% Punctuation after theorem head
	{.5em}	% Space after theorem head
	{\thmname{#1}\thmnumber{ #2}\thmnote{ \normalfont{\text{(#3)}}}}	% head content
\newtheoremstyle{c_definition}
	{3pt}	% Space above
	{3pt}	% Space below
	{}% Body font
	{}	% Indent amount
	{\bfseries}	% Theorem head font
	{}	% Punctuation after theorem head
	{.5em}	% Space after theorem head
	{\thmname{#1}\thmnumber{ #2}\thmnote{ \normalfont{\text{(#3)}}}}	% head content
\newtheoremstyle{c_plain}
	{3pt}	% Space above
	{3pt}	% Space below
	{\itshape}% Body font
	{}	% Indent amount
	{\bfseries}	% Theorem head font
	{}	% Punctuation after theorem head
	{.5em}	% Space after theorem head
	{\thmname{#1}\thmnumber{ #2}\thmnote{ \text{(#3)}}}	% head content

\ifcsname c@english\endcsname
	\theoremstyle{plain}
	\newtheorem{theorem}{Theorem}[section]
	\newtheorem*{theorem*}{Theorem}
	\newtheorem{lemma}[theorem]{Lemma}
	\newtheorem{proposition}[theorem]{Proposition}
	\newtheorem*{proposition*}{Proposition}

	\theoremstyle{definition}
	\newtheorem{definition}[theorem]{Definition}
	\newtheorem*{definition*}{Definition}
	\newtheorem{example}{Example}[section]
	\newtheorem{exercise}{Exercise}[section]

	\theoremstyle{remark}
	\newtheorem*{remark}{Remark}
	\newtheorem*{solution}{Solution}
	\newtheorem{corollary}[theorem]{Corollary}
	\newtheorem{notation}[theorem]{Notation}
\else
	\theoremstyle{c_plain}
	\newtheorem{theorem}{משפט}[section]
	\newtheorem*{theorem*}{משפט}
	\newtheorem{lemma}[theorem]{למה}
	\newtheorem{proposition}[theorem]{טענה}
	\newtheorem*{proposition*}{טענה}

	\theoremstyle{c_definition}
	\newtheorem{definition}[theorem]{הגדרה}
	\newtheorem*{definition*}{הגדרה}
	\newtheorem{example}{דוגמה}[section]
	\newtheorem{exercise}{תרגיל}[section]

	\theoremstyle{c_remark}
	\newtheorem*{remark}{הערה}
	\newtheorem*{solution}{פתרון}
	\newtheorem{corollary}[theorem]{מסקנה}
	\newtheorem{notation}[theorem]{סימון}
\fi

% Questions related commands
\newcounter{question}
\setcounter{question}{1}
\newcounter{sub_question}
\setcounter{sub_question}{1}

\ifcsname c@english\endcsname
	\newcommand{\question}[1][0]{
		\ifthenelse{#1 = 0}{}{\setcounter{question}{#1}}
		\section{Question \arabic{question}}
		\addtocounter{question}{1}
		\setcounter{sub_question}{1}
	}

	\newcommand{\subquestion}[1][0]{
		\ifthenelse{#1 = 0}{}{\setcounter{sub_question}{#1}}
		\subsection{Part \alph{sub_question}}
		\addtocounter{sub_question}{1}
	}
\else
	\newcommand{\question}[1][0]{
		\ifthenelse{#1 = 0}{}{\setcounter{question}{#1}}
		\section{שאלה \arabic{question}}
		\addtocounter{question}{1}
		\setcounter{sub_question}{1}
	}

	\newcommand{\subquestion}[1][0]{
		\ifthenelse{#1 = 0}{}{\setcounter{sub_question}{#1}}
		\subsection{סעיף \localecounter{letters.gershayim}{sub_question}}
		\addtocounter{sub_question}{1}
	}
\fi

% import lua and start of document
\directlua{common = require ('../common')}

\GetEnv{AUTHOR}

% headers
\author{\AUTHOR}
\date\today

\setcounter{secnumdepth}{2}
% chktex-file 9
% chktex-file 17

\title{מתמטיקה א' --- תרגול 2}
\date{}

\begin{document}

\maketitle

\section{קבוצות וסימונים}
\begin{definition}[קבוצה]
	קבוצה היא אוסף של איברים. בקבוצה אין חשיבות לסדר האיברים ולכפילות של הופעת איבר. קבוצה של איברים תסומן באמצעות סוגריים מסולסלים $\{\}$ ואיבריה יופרדו באמצעות פסיק.
\end{definition}
\begin{remark}
	בקורס שלנו נעסוק כמעט תמיד בקבוצות של מספרים. את הקבוצות נסמן לרוב באמצעות אות אנגלית גדולה. \\
	לדוגמה, נוכל להגדיר את הקבוצה $A = \{1, 2, 3\}$. מכיוון שאין חשיבות לסדר ולכפילויות:
	\[
		A = \{1, 2, 3\} = \{2, 1, 3\} = \{1, 1, 2, 3, 2\}
	\]
\end{remark}
פעמים רבות נרצה להעביר רעיון מתמטי באמצעות סימנים ולא באמצעות מילים. כך למשל כולנו מבינים שהמשמעות של $3 + 5$ היא שלוש ועוד חמש, ואנחנו לא נדרשים לכתוב את המילים.

נכיר כעת מספר סימנים מתמטיים בהם נשתמש רבות:
\begin{itemize}
    \item $\in$ --- סימן השייכות. נכתוב $a \in A$ אם נרצה לומר ''האיבר $a$ שייך לקבוצה $A$''.
    \item $\notin$ --- סימן אי השייכות. נכתוב $a \notin A$ אם נרצה לומר ''האיבר $a$ איננו שייך לקבוצה $A$''.
    \item $\forall$ --- סימן ''לכל''.
    \item $\exists$ --- סימן ''קיים''.
    \item $\nexists$ --- סימן ''לא קיים''.
    \item $\exists! $ --- סימן ''קיים ויחיד''.
\end{itemize}
\begin{example}
	כדי לכתוב שלכל מספר טבעי יש מספר נוסף הגדול ממנו נוכל לכתוב,
	\[
		\forall x \in \{1, 2, 3, \ldots\} \exists y \ y > x
	\]
	כדי לכתוב שלא קיים מספר הגדול מכל מספר טבעי נכתוב,
	\[
		\nexists y \forall x \in \{1, 2, 3, \ldots\}\ y > x
	\]
\end{example}
\begin{itemize}
    \item $\subseteq$ --- סימן ההכלה. נגיד שהקבוצה $A$ מוכלת בקבוצה $B$, ונכתוב $A \subseteq B$, אם כל איברי הקבוצה $A$ הם גם איברים של הקבוצה $B$. \\
    למשל עבור $A = \{1, 2, 3\}$ ו־$B = \{1, 2, 3, 4\}$ מתקיים $A \subseteq B$.
    \item $\nsubseteq$ --- סימן אי הכלה. נגיד שהקבוצה $A$ לא מוכלת בקבוצה $B$, ונסמן $A \nsubseteq B$ אם יש ב־$A$ איברים שאינם איברים של הקבוצה $B$. \\
    למשל עבור $A = \{1, 2, 3\}$ ו־$B = \{1, 2, 3, 4\}$ מתקיים $B \nsubseteq A$.
    \item שוויון קבוצות. נגיד $A = B$ אם בשתי הקבוצות יש בדיוק אותם איברים.
    \item $\emptyset = \{\}$ --- הקבוצה הריקה, כלומר קבוצה ללא איברים.
\end{itemize}
\begin{remark}
	הקבוצה הריקה מוכלת בכל קבוצה. מדוע? זאת מכיוון שלא קיימים בקבוצה הריקה איברים שלא שייכים לקבוצה אחרת.
\end{remark}
\begin{example}
	נתבונן בקבוצה $A = \{1, \{1\}, \{2\}\}$.
	\begin{itemize}
		\item האם $1 \in A$? כן.
		\item האם $2 \in A$? לא.
		\item האם $\{2\} \in A$? כן.
		\item האם $\{\{1\}, \{2\}\} = A$? לא.
		\item האם $\{\{1\}, \{2\}\} \subseteq A$? כן.
	\end{itemize}
\end{example}

\section{דרכים לבנות קבוצות}
\subsection{הגדרת קבוצה}
ישנן שלוש דרכים בהן נגדיר קבוצות לאורך הקורס. בעיקר חשוב לוודא שאנחנו יודעים כיצד לקרוא הגדרה של קבוצה.
\begin{enumerate}
    \item \textbf{הגדרה מפורשת.} למשל $A = \{1, 2, 3\}$ או $C = \{1, 3, 5, 7, 9, \ldots\}$.
    \item \textbf{עיקרון ההפרדה.} נשתמש בו כאשר נרצה להפריד את הקבוצה $A$ מתוך קבוצה $B$ קיימת. למשל, אם $B = \{1, 2, 3, 4\}$, נוכל להגדיר את הקבוצה $A = \{1, 2, 3\}$ בצורה הבאה:
    \[
        A = \{x \in B \mid x < 4\}
    \]
    \item \textbf{עיקרון ההחלפה.} נשתמש בעיקרון זה כשנרצה להחליף איברים בקבוצה מסוימת. למשל: את הקבוצה $D = \{2, 4, 6\}$ נוכל להגדיר בצורה הבאה:
    \[
        D = \{2 \cdot x \mid x \in A\}
    \]
\end{enumerate}

\subsection{קבוצות המספרים}
ישנן מספר קבוצות מספרים חשובות שנשתמש בהן לאורך הקורס פעמים רבות, ולכן נקבע להן סימון מיוחד.
\begin{itemize}
    \item $\NN = \{1, 2, 3, \ldots\}$ --- קבוצת המספרים הטבעיים
    \item $\ZZ = \{0, -1, 1, -2, 2, -3, 3, \ldots\}$ --- קבוצת המספרים השלמים
    \item $\QQ = \left\{ \frac{m}{n} \;\middle|\; m, n \in \ZZ, n \neq 0 \right\}$ --- קבוצת המספרים הרציונליים
    \item $\RR = \{x \mid -\infty < x < \infty\}$ --- קבוצת המספרים הממשיים
\end{itemize}
נבחין כי מתקיים:
\[
    \NN \subseteq \ZZ \subseteq \QQ \subseteq \RR
\]

\subsection{קטעים וקרנות}
סוג נוסף של קבוצות בהן נתקל רבות בקורס הוא קטע וקרן. \\
בהינתן $a, b \in \RR$, נגדיר:
\begin{enumerate}
    \item \textbf{קטע סגור (כולל הקצוות):}
    \[
        [a, b] = \{x \in \RR \mid a \le x \le b\}
    \]
    \item \textbf{קטע פתוח (לא כולל הקצוות):}
    \[
        (a, b) = \{x \in \RR \mid a < x < b\}
    \]
    \item \textbf{קטע חצי סגור מצד שמאל (כולל הקצה השמאלי):}
    \[
        [a, b) = \{x \in \RR \mid a \le x < b\}
    \]
    \item \textbf{קטע חצי סגור מצד ימין (כולל הקצה הימני):}
    \[
        (a, b] = \{x \in \RR \mid a < x \le b\}
    \]
    \item \textbf{קרן סגורה משמאל:}
    \[
        [a, \infty) = \{x \in \RR \mid a \le x\}
    \]
    \item \textbf{קרן פתוחה משמאל:}
    \[
        (a, \infty) = \{x \in \RR \mid a < x\}
    \]
    \item \textbf{קרן סגורה מימין:}
    \[
        (-\infty, a] = \{x \in \RR \mid x \le a\}
    \]
    \item \textbf{קרן פתוחה מימין:}
    \[
        (-\infty, a) = \{x \in \RR \mid x < a\}
    \]
\end{enumerate}
\begin{remark}
	נשים לב כי $[a, a] = \{a\}$, וכן $(a, a) = \emptyset$. בנוסף, אם $a > b$, הקטע $[a, b] = \emptyset$. \\
	כמו כן, מכיוון ש־$\infty$ איננו מספר, אין משמעות לסימון $[a, \infty]$.
\end{remark}
\subsection{איחוד וחיתוך}
פעולות האיחוד והחיתוך מאפשרות לנו לייצר קבוצה חדשה מתוך שתי קבוצות קיימות. בהינתן $A, B$ קבוצות:
\begin{definition}[איחוד וחיתוך]
	הקבוצה $A \cup B$ לה נקרא ''$A$ איחוד $B$'' היא הקבוצה:
	\[
		A \cup B = \{x \mid x \in A \text{ או } x \in B\}
	\]
	הקבוצה $A \cap B$ לה נקרא ''$A$ חיתוך $B$'' היא הקבוצה:
	\[
		A \cap B = \{x \mid x \in A \text{ וגם } x \in B\}
	\]
\end{definition}
\begin{example}
	עבור $A = [1, 2)$ ו־$B = (1.5, 3]$,
	\[
		A \cup B = [1, 3]
	\]
	וכן,
	\[
		A \cap B = (1.5, 2)
	\]
\end{example}
\begin{remark}
	נשים לב שלעתים אין דרך ''פשוטה יותר'' לתאר קבוצת איחוד או חיתוך. למשל, הדרך הנוחה ביותר לתאר את הקבוצה $[0, 1] \cup [2, 3]$ היא $[0, 1] \cup [2, 3]$.
\end{remark}
\subsection{דוגמאות נוספות}
חשבו את הקבוצות הבאות:
\begin{itemize}
    \item $\{|x| \mid x \in \ZZ, x \neq 0\} = \NN$
    \item $\{x - 1 \mid x \in \NN, x \ge 2\} = \NN$
    \item $\{2x \mid x \in \NN\} = \{2, 4, 6, \ldots\}$
    \item $\{x^2 \mid x \in \RR\} = [0, \infty)$
    \item $\NN \cup \{0\} \cup \{-x \mid x \in \NN\} = \ZZ$
    \item $[0, 4) \cap [400, \infty) = \emptyset$
    \item $(0, 7) \cup \{0, 3, 6\} = [0, 7)$
    \item $(-\infty, 6) \cup (3, \infty) = \RR$
\end{itemize}

\subsection{פונקציות}
\begin{definition}[פונקציה]
	פונקציה $f$ היא מיפוי בין שתי קבוצות $A$ ו־$B$ המקיים את כלל הפונקציה: \\
	לכל $x \in A$ קיים איבר יחיד $y \in B$ אשר הפונקציה מקבלת מ־$x$, נסמן זאת $f(x) = y$ ובמקרה זה נסמן $f : A \to B$.
\end{definition}
\begin{example}
	אם $A = \{1, 2, 3\}$ ו־$B = \{4, 5\}$ אז לדוגמה נוכל להגדיר $f(1) = f(2) = 4, f(3) = 5$.
\end{example}
\begin{definition}[תחום, טווח ותמונה]
	בהינתן $f : A \to B$ פונקציה נאמר ש־$A$ הוא התחום של $f$, ו־$B$ הוא הטווח.
	בנוסף ה\textbf{תמונה} של $f$ היא קבוצת כל הערכים ש־$f$ יכולה לקבל, נסמן $\im f \subseteq B$.
\end{definition}
\begin{example}
	נגדיר את $f(x) = x^2$, זוהי פרבולה (מחייכת).
	מה התחום של הפונקציה? $\RR$ (קבוצת הממשיים). \\
	מה הטווח שלה? למעשה אנו יכולים להחליט זאת, אם נסמן את הטווח $\QQ$ (הרציונליים) לדוגמה, אז $f$ לא תהיה פונקציה תקינה, שהרי התוצאה של $f(\sqrt[3]{2}) \notin \QQ$, לכן לדוגמה נקבל את הטווח להיות $\RR$. \\
	מה התמונה של $f$? ריבוע תמיד מחזיר ערכים חיוביים או אפס (מספרים אי־שליליים) ולכן במקרה זה $\im f = \{ x \ge 0 \mid x \in \RR \}$.
\end{example}
\begin{exercise}
	מצאו את התחום הטווח ותמונה של $f : A \to B$ בסעיפים הבאים:
	\begin{enumerate}
		\item $f(x) = 2x - 3$
		\item $f(x) = \frac{1}{x - 1}$
		\item $f(x) = \sqrt{x}$
		\item $f(x) = 0$ אם $x$ הוא מספר זוגי ו־$f(x) = 1$ אם היא מספר אי־זוגי
	\end{enumerate}
\end{exercise}
\begin{solution}
	\begin{enumerate}
		\item $A = \RR, B = \RR, \im = \RR$, אבל גם $A = [0, 1], B = [-3, -1], \im = B$ יעבוד
		\item $A = \RR \setminus \{ 1 \}$ שכן $f(1) = \frac{1}{1 - 1}$ לא מוגדר, $B = \RR$ ו־$\im f = \RR \setminus \{ 0 \}$
		\item $A = \{ x \ge 0 \mid x \in \RR \}, B = \RR, \im = A$
		\item במקרה זה הפונקציה מוגדרת רק עבור מספרים טבעיים (זוגיים ואי־זוגיים שהם כל הטבעיים) ולכן $A = \NN$, הטווח חייב להכיל את $0, 1$ ולכן נקבע $B = \RR$ וכן $\im f = \{ 0, 1 \}$
	\end{enumerate}
\end{solution}

\end{document}
